% TOP_PROBLEM: {"id": 3318, "title": "The Crossing Number of the Complete Graph", "category_id": 3, "category": {"id": 3, "name": "graph_theory", "display_name": "Graph Theory", "description": "Problems involving graphs, networks, and their properties.", "slug": "graph-theory", "order_index": 3, "created_at": "2026-07-31T15:26:25.671Z"}, "status": "open", "problem_number": "OPG-307", "set_id": 12}
% FIRST_POSTED: 2026-10-01
% ENTRY_KIND: statement_only
% UnsolvedMath ID 3318; source catalogue code OPG-307
% !TeX program = lualatex
% Definition and sources only; this file is not an LLM solution attempt.
\documentclass[11pt,a4paper]{article}
\usepackage[margin=25mm]{geometry}
\usepackage{fontspec}
\setmainfont{lmroman10-regular.otf}[BoldFont=lmroman10-bold.otf,
 ItalicFont=lmroman10-italic.otf,BoldItalicFont=lmroman10-bolditalic.otf]
\usepackage{amsmath,amssymb,mathtools}
\usepackage{unicode-math}
\setmathfont{latinmodern-math.otf}
\AtBeginDocument{\let\setminus\smallsetminus}
\usepackage{xurl,hyperref}
\hypersetup{colorlinks=true,urlcolor=blue}
\setcounter{secnumdepth}{0}
\setlength{\parindent}{0pt}
\setlength{\parskip}{5pt}
\setlength{\emergencystretch}{3em}
\begin{document}
\section*{The Crossing Number of the Complete Graph}\label{3318}
{\small UnsolvedMath ID 3318 · OPG-307 · Graph Theory · OpenGarden}
\subsection{Definitions and mathematical statement}
Let $K_n$ be the finite simple complete graph on $n\ge1$ vertices.
A topological drawing in the plane represents vertices by distinct
points and edges by simple arcs between their endpoints. No edge
passes through another vertex. Edge interiors meet, if at all, in
finitely many transverse crossings, with no three edges crossing
at one point. Common endpoints are not crossings. The crossing
number $\operatorname{cr}(G)$ is the minimum total number of such
crossings among all drawings of a graph $G$.

The conjecture asks whether, for every integer $n\ge1$,
\[
 \operatorname{cr}(K_n)=
 \frac14\left\lfloor\frac n2\right\rfloor
          \left\lfloor\frac{n-1}2\right\rfloor
          \left\lfloor\frac{n-2}2\right\rfloor
          \left\lfloor\frac{n-3}2\right\rfloor.
\]
The floor symbols mean rounding down. This expression is zero for
$n\le4$, consistent with planar drawings, and gives one for $n=5$.
The source records constructions attaining the displayed expression;
the asserted equality requires that no drawing do better.

Edges may bend freely: no requirement of straight-line segments,
convex vertex placement, or a fixed cyclic ordering of vertices is
imposed. Thus the problem concerns topological crossing number,
which differs from the rectilinear parameter. Crossings are counted
as points, so if a drawing lets the same two edges cross more than
once, every crossing contributes. A proof of equality must cover
all finite orders, rather than only showing that the ratio to the
displayed expression tends to one.
\subsection{Short English statement}
Is the crossing number of the complete graph exactly the displayed four-factor floor formula for every number of vertices?
\subsection{Sources}
\begin{itemize}
\item[S1] ULAM.AI and the UnsolvedMath Contributors, supplied problems.json, record 3318 (OPG-307), OpenGarden collection. Catalogue identity and source transcription; CC BY 4.0.
\url{https://huggingface.co/datasets/ulamai/UnsolvedMath}.
\item[S2] Open Problem Garden, The Crossing Number of the Complete Graph. Original problem and discussion; the mathematical formulation below is expanded from the supplied source record.
\url{http://www.openproblemgarden.org/op/the_crossing_number_of_the_complete_graph}.
\end{itemize}
\end{document}
